%% kapa: exam=1 round=23 year=2012 period=2 form=A subject=appraisal_law subject_ko=부동산관계법규 seq=26 number=66
%% source: Q-Net 감정평가사 기출문제 (cst00309) artlSeq=5052580 file=raw/1008/5052580/2012년 제23회 감정평가사 제1차시험 기출문제(2교시 A형).pdf
%% conversion: kordoc 4.24.51 -> tools/gen1.py (auto); review=proofread-sonnet; flags=none
\begin{question}{66}
국유재산법령상 국유재산의 종류와 범위 등에 관한 설명으로 옳지 않은 것은?
\begin{choices}
  \item 정부기업이 직원의 주거용으로 사용하는 국유재산은 일반재산이다.
  \item 국유재산을 관리ㆍ처분할 경우에는 취득과 처분이 균형을 이루어야 한다.
  \item 디자인권은 국유재산의 범위에 포함된다.
  \item 국유재산은 행정재산과 일반재산으로 구분되며, 행정재산은 공용재산, 공공용재산, 기업용재산, 보존용재산으로 분류된다.
  \item 정부기업이 사용하던 궤도차량이 정부기업의 폐지와 함께 포괄적으로 용도폐지된 경우에도 국유재산에 해당한다.
\end{choices}
\end{question}
\answer{1}
